\section{Introduction}
Let $\Omega \Subset \C^n$ be a bounded domain and $m$ an integer such that $1 \leq m \leq n$. We~deno\-te by $\mathcal{SH}_m (\Omega)$ the set of $m$-subharmonic functions on $\Omega$. The Dirichlet problem we are concerned with is the following: find $U \in \mathcal{SH}_m (\Omega) \cap C^0 (\bar{\Omega})$ such that
\begin{equation}\label{eq:DirPb}
\begin{cases}
(dd^c U)^m \wedge \beta^{n - m} = \mu &\text{on}\ \Omega, \\
U_{\mid \partial \Omega} = g & \text{on}\ \partial \Omega,
\end{cases}
\end{equation}
where the boundary datum $g \in C^0 (\partial \Omega)$ and $\mu$ is a positive Borel measure on $\Omega$ with finite mass $\mu (\Omega) < + \infty$.

The Hölder continuous subsolution problem asks basically if the Dirichlet problem \eqref{eq:DirPb} admits a Hölder continuous solution when it admits a Hölder continuous subsolution, provided that the boundary datum $g$ is Hölder continuous.

This qualitative statement is true and was essentially proved in \cite[Th.\,B]{BZ20} and is essentially correct when the boundary datum $g$ belongs to $C^{1,1} (\partial \Omega)$. However its proof in the general case relies on \cite[Lem.\,4.2]{BZ20} whose proof is not complete (see Lemma \ref{lem:ModC} below).

Here is the correct version of Theorem B of \cite{BZ20}.

\begin{theoremB}\label{th:B'}
Let $\Omega \Subset \C^n$ be a bounded strongly $m$-pseudoconvex domain and $\mu $ a positive Borel measure on $\Omega$ with finite mass. Assume that there exists \hbox{$\varphi\in \mathcal{E}^0_m(\Omega)\cap\mathcal{C}^{\alpha}(\overline\Omega)$} with $0 < \alpha < 1$ such that
\begin{equation} \label{eq:subsol}
\mu \leq (dd^c\varphi)^m\wedge\beta^{n-m}, \quad \text{weakly on $\Omega$, and $\varphi|_{\partial \Omega} \equiv 0$.}
\end{equation}

Then for any boundary datum $g $ Hölder continuous on $\partial \Omega$, the Dirichlet problem \eqref{eq:DirPb} admits a unique solution $U = U_{g,\mu} $ which is Hölder continuous on $\bar{\Omega}$.
More precisely,
\begin{enumerate}
\item
if $g \in \mathcal C^{1,1} (\partial \Omega)$, then $ U \in \mathcal{C}^{\alpha'}(\overline\Omega)$ for any $\alpha'$ such that
\[
0 < \alpha' < 2 \gamma' (m,n,\alpha) \sfrac{\alpha^m}{2^{m}},
\]
where
\begin{equation} \label{eq:gamma'}
\gamma' (m,n,\alpha) := \frac{ m \alpha}{ m (m + 1) \alpha + (n-m) [(2 - \alpha)m + \alpha]},
\end{equation}

\item
if $g \in \mathcal C^{2 \alpha} (\partial \Omega)$, then $ U \in \mathcal{C}^{\alpha''}(\overline\Omega)$ for any $\alpha''$ such that
\[
0 < \alpha'' < \gamma'' (m,n,\alpha) \sfrac{\alpha^m}{2^{m}},
\]
where
\begin{equation} \label{eq:gamma''}
\gamma'' (m,n,\alpha) := \frac{\alpha}{ m (m + 1) \alpha + (n-m) [(2 - \alpha)m + \alpha]}\cdot
\end{equation}
\end{enumerate}
\end{theoremB}

The first statement of Theorem \ref{th:B'} is the same as that of \cite[Th.\,B]{BZ20} when the boundary datum $g \in \mathcal C^{1,1} (\partial \Omega)$, except for the exponent.
In fact the one given in the statement of \cite[Th.\,B]{BZ20} is not correct due to an unfortunate misprint at the end of its proof in page 1005, where \cite[Prop.\,2.20]{BZ20} was applied with a wrong exponent.

On the other hand, when $g$ is only assumed to be Hölder continuous on $\partial \Omega$, we had to find a new argument because in this case we do not know if the bounds given by \eqref{eq:claim2} in Lemma \ref{lem:claim2} below still hold. This will be done in Lemma \ref{lem:ModC} below, but unfortunately this leads to a worse exponent.

\subsubsection*{Acknowledgements}
We thank the referee for his careful reading of this note and his useful suggestions that helped to improve the presentation of the mains results.

\section{A new version of \texorpdfstring{\cite[Lem.\,4.2]{BZ20}}{BZ20}}

Here we will give a new version of \cite[Lem.\,4.2]{BZ20} and its complete proof following essentially the same scheme. Before stating this new version, we need to prove the following result which provides the argument missing in the proof of \cite[Lem.\,4.2]{BZ20}.

\begin{lemma} \label{lem:claim2}
Let $\Omega \Subset \C^n$ be a bounded strongly $m$-pseudoconvex domain and $k$ an integer such that $0 \leq k \leq m-1$.
\begin{enumerate}
\item
Assume that $\psi \in \mathcal{SH}_m (\Omega) \cap L^{\infty} (\Omega)$, $u,v\in \mathcal{SH}_m(\Omega)\cap L^{\infty}(\Omega)$ satisfy $u \leq v$ on~$\Omega$ and for any $\zeta \in \partial \Omega$, $\lim_{z \rightarrow \zeta }(u(z)- v(z))= 0$.
Then
\begin{equation} \label{eq:claim1}
\int_{\Omega} dd^c v \wedge (dd^c \psi)^{k} \wedge\beta^{n-k-1} \leq \int_{\Omega} dd^c u \wedge (dd^c \psi)^k \wedge\beta^{n-k-1}.
\end{equation}

\item
Assume that $g \in C^{1,1} (\partial \Omega)$. Then there exists a constant $ M' = M'(m,n,g) > 0$ such that for any $\psi \in \mathcal{E}^0_m (\Omega)$ and $v \in \mathcal{SH}_m (\Omega,R)$ with $v|_{\partial \Omega} = g$, we have
\begin{equation} \label{eq:claim2}
\int_{\Omega} dd^c v \wedge(dd^c \psi)^k\wedge\beta^{n-k-1} \leq (R + M') \, H_m(\psi)^{k\slash m} .
\end{equation}
\end{enumerate}
\end{lemma}
The first statement is well known and was stated in \cite[Cor.\,2.9]{BZ20}. Since its proof there was not complete, we will give a different proof here.

\skpt
\begin{proof}
\begin{enumerate}
\item
Fix $\varepsilon > 0$. From the hypothesis, the exists a compact subset $K \Subset \Omega$ such that $ u \geq v - \varepsilon$ on $\Omega \setminus K$. Then $v_\varepsilon := \max \{u,v-\varepsilon\} \in \mathcal{SH}_m (\Omega) \cap L^{\infty} (\Omega)$ and $v_\varepsilon = u$ on $\Omega \setminus K$.
We claim that
\begin{equation} \label{eq:masspreservation}
\int_{\Omega} dd^c v_\varepsilon \wedge (dd^c \psi)^{k} \wedge\beta^{n-k-1} = \int_{\Omega} dd^c u \wedge (dd^c \psi)^k \wedge\beta^{n-k-1}.
\end{equation}
Indeed we have in the sense of currents on $\Omega$,
\[
dd^c v_\varepsilon \wedge (dd^c \psi)^{k} \wedge\beta^{n-k-1} - dd^c u \wedge (dd^c \psi)^k \wedge\beta^{n-k-1} = dd^c T,
\]
where $T := (v_\varepsilon - u) (dd^c \psi)^k \wedge\beta^{n-k-1}$.

Since $T$ is a current of order $0$ with compact support in $\Omega$, it follows that $\int_\Omega dd^c T =\nobreak 0$, which proves \eqref{eq:masspreservation}.
Now observe that $(v_\varepsilon)$ increases to $v$ as $\varepsilon$ decreases to $0$. Therefore, taking the limit in \eqref{eq:masspreservation} as $\varepsilon \to 0$ and using the monotone continuity of the Hessian operators, we obtain the inequality \eqref{eq:claim1}.

\item
Let us first assume that $g|_{\partial \Omega} \equiv 0$. Then $v|_{\partial \Omega} \equiv 0$ and we can apply Cegrell type inequalities \cite[Lem.\,2.8]{BZ20} and use the normalization mass condition, to conclude that
\[
I_k(v,\psi) \leq H_m (v)^{1\slash m} H_m (\psi)^{k \slash m}
\leq  R \, H_m (\psi)^{k \slash m}.
\]
This proves the inequality \eqref{eq:claim2}.

Now assume that $g \in C^{1,1} (\partial \Omega)$ and fix $k$ such that $1 \leq k \leq m$.
There exists $G \in C^{1,1} (\bar{\Omega})$ such that $G = g$ on $\partial \Omega$. Since $\Omega$ is strongly $m$-pseudoconvex, it admits a defining function $\rho$ which is strongly $m$-subharmonic on $\Omega$. Then we can find a large constant $L > 0$ such that $w := L \rho + G$ is $m$-subharmonic on $\Omega$ and for $1 \leq k \leq m$, $(dd^c w)^k \wedge \beta^{n-k} \leq L'_m \beta^n$ pointwise almost everywhere in $\Omega$ for some uniform constant $L'_m > 0$.

By the bounded subsolution theorem (\cite{N13}) there exists $v_0 \in \mathcal{E}^0_m (\Omega)$ solution to the equation $(dd^c v_0)^m \wedge \beta^{n-m} = (dd^c v)^m \wedge \beta^{n-m}$ with boundary values $v_0 \equiv 0$.
The function $\tilde v := v_0 + w$ is $m$-subharmonic and bounded on $\Omega$, $(dd^c \tilde v)^m \wedge \beta^{n-m} \geq (dd^c v)^m \wedge \beta^{n-m} $ on $\Omega$ and $ \tilde v = g = v$ on $\partial \Omega$. By the comparison principle we get $\tilde v \leq v$ on $\Omega$.
Therefore, by the formula \eqref{eq:claim1} we conclude that
\[
I_k (v,\psi) \leq I_k (\tilde v,\psi).
\]

It suffices to estimate $ I_k (\tilde v,\psi) $ by a uniform constant. We have
\[
I_k (\tilde v,\psi) = I_k (v_0, \psi) + I_k (w,\psi).
\]
Since $v_0|_{\partial \Omega} \equiv 0$, from the previous case it follows that
\begin{align*}
I_k (v_0,\psi) & \leq  \left(\int_\Omega (dd^c v_0)^m \wedge \beta^{n-m}\right)^{1 \slash m} \left(\int_\Omega (dd^c \psi)^m \wedge \beta^{n-m}\right)^{k \slash m} \\
&\leq  H_m (\psi)^{k\slash m} R.
\end{align*}

It remains to estimate $ I_k (w,\psi)$.
Since $w \in C^{1,1} (\bar{\Omega})$, it follows that $dd^c w \leq M_3 \beta$ pointwise almost everywhere on $\Omega$, hence by \cite[Lem.\,2.8]{BZ20} , we have
\begin{align*}
\int_\Omega dd^c w \wedge (dd^c \psi)^k \wedge \beta^{n-k-1} &\leq  M' \int_\Omega (dd^c \psi) ^k \wedge \beta^{n-k} \\
& \leq  M' H_m (\psi)^{k\slash m} \leq M' H_m (\psi)^{k\slash m},
\end{align*}
where $M' = M'(g) > 0$ depends on the uniform bound of $dd^c G$.\qedhere
\end{enumerate}
\end{proof}

To state the main lemma which will replace \cite[Lem.\,4.2]{BZ20}, let us fix some notations.

Fix $R \geq 1$ and denote by $\mathcal{SH}_m (\Omega,R)$) the set of bounded negative $m$-subharmonic functions $w$ on $\Omega$ such that $\int_\Omega (dd^c w)^m \wedge \beta^{n-m} \leq R^m$.
On the other hand, we denote by $ \mathcal{E}_m^0 (\Omega,R) = \mathcal{SH}_m (\Omega,R) \cap \mathcal{E}^0_m (\Omega)$.

\begin{lemma}\label{lem:ModC} Let $\Omega \Subset \C^n$ be a bounded strongly $m$-pseudoconvex domain and $\varphi\in \mathcal{E}^0_m (\Omega)\cap \mathcal{C}^{\alpha} (\overline\Omega)$, with $ 0 < \alpha \leq 1$. Then for every $k$ such that $1 \leq k \leq m$, there exists a constant $\tilde C_k = \tilde C (k,m,\varphi, \Omega) >0$ such that for every $u,v\in \mathcal{SH}_m (\Omega,R)$ such that $u = v$ on $\partial \Omega$, we have
\begin{equation} \label{eq:MocEst2}
\int_{\Omega}|u-v| (dd^c\varphi)^k\wedge\beta^{n-k}\leq \tilde C_k (1 + \Vert u - v\Vert_{\infty}^{m-1}) \, R \, \left[\Vert u-v \Vert_1\right]^{\tilde \alpha_k },
\end{equation}
provided that $\Vert u-v \Vert_1 \leq 1$, where $\tilde \alpha_k := \sfrac{\alpha^k}{m 2^k} $.

Furthermore if $g \in C^{1,1} (\partial \Omega)$, for any $k$ such that $1 \leq k\leq m$, there exists a constant $C_k = C (k,m, \varphi,g,\Omega) > 0$ such that for every $u, v \in \mathcal{SH}_m (\Omega,R)$ with $u = g = v$ on $\partial \Omega$, we~have
\begin{equation} \label{eq:MocEst1}
\int_{\Omega}|u-v|(dd^c\varphi)^k\wedge\beta^{n-k}\leq C_k R \left[\Vert u-v \Vert_1\right]^{\alpha_k},
\end{equation}
provided that $\Vert u-v \Vert_1 := \int_{\Omega} \vert u - v\vert \, \beta^n \leq 1$, where $\alpha_k := 2 (\sfrac{\alpha}{ 2})^k$.
\end{lemma}

The proof follows the same scheme as that of \cite[Lem.\,4.2]{BZ20}. More precisely,
we will first start by assuming $g|_{\partial \Omega} \equiv 0$, since the proof is exactly the same in this case.
Then we will consider the case $g\in C^{1,1} (\partial \Omega)$, which can be reduced to the previous case using Lemma \ref{lem:claim2}.
Finally we will treat the general case by emphasizing the new ingredient to be used to fill the gap in the proof of \cite[Lem.\,4.2]{BZ20}.

\begin{proof}
Let us first recall the main steps of the proof of \cite[Lem.\,4.2]{BZ20}.
We can reduce to the case when $u = v$ near the boundary $\partial \Omega$ and $u \geq v$ on $\Omega$.

We approximate $\varphi$ by smooth functions. As in \cite{BZ20}, we extend $\varphi$ as a Hölder continuous function of order $\alpha$ on $\C^n$ and denote by $\varphi_{\delta}$ ($0 < \delta < \delta_0$) the usual smooth approximants of $\varphi $ on $\C^n$.
We know that $\varphi_\delta \in \mathcal{SH}_m (\Omega_{\delta})\cap\mathcal{C}^{\infty}(\mathbb C^n)$.

To prove the required estimates, we will argue by induction on $k$ such that \hbox{$0\!\leq\!k\!\leq\!m$}. Fix $0 \leq k \leq m-1$ and $\delta_0 > 0$ small enough and write for $\delta$ such that $0 < \delta < \delta_0$,
\[
\int_{\Omega}(u-v) (dd^c\varphi)^{k+1} \wedge\beta^{n-k-1} = A (\delta) + B (\delta),
\]
where
\[
A (\delta):=\int_{\Omega}(u-v)dd^c \varphi_{\delta}\wedge(dd^c\varphi)^k\wedge\beta^{n-k-1},
\]
and
\[
B (\delta) := \int_{\Omega}(u-v) dd^c(\varphi - \varphi_{\delta} - \kappa \delta^\alpha)\wedge(dd^c\varphi)^k\wedge\beta^{n-k-1},
\]
where we recall that $ \varphi - \kappa \delta^\alpha \leq \varphi_\delta \leq \varphi + \kappa \delta^\alpha$ on $\bar{\Omega}$ by hypothesis.

The first term $A (\delta)$ is estimated as follows. Observe that we have
\[
dd^c \varphi_\delta \leq M_1 \kappa \delta^{\alpha -2} \beta
\]
pointwise on $\Omega$, where $M_1 > 0$ is a uniform bound on the second derivatives of $\chi$. Then since $u \geq v$ we deduce that
\begin{equation} \label{eq:estimationA}
\vert A (\delta) \vert \leq M_1 \kappa \delta^{\alpha -2} \int_{\Omega} (u-v) (dd^c \varphi)^k \wedge \beta^{n-k}.
\end{equation}

We now estimate the second term $B (\delta)$ following the arguments of \cite[Lem.\,4.2]{BZ20} with a slight modification.

Since $ u - v = 0$ near the boundary $\partial \Omega$, \ie on $\Omega \setminus \Omega'$, where $\Omega' \Subset \Omega$ is an open set, we can integrate by parts to get the following formula
\[
B (\delta) = \int_{\Omega'} (\varphi_{\delta} - \varphi + \kappa \delta^\alpha) dd^c (v-u) \wedge(dd^c\varphi)^k\wedge\beta^{n-k-1},
\]
and then, since $0 \leq \varphi_{\delta}-\varphi + \kappa \delta^\alpha \leq 2 \kappa \delta^{\alpha}$ on $\Omega$, it follows that
\[
\vert B (\delta)\vert \leq 2 \kappa \delta^{\alpha} \int_{\Omega'} dd^c v \wedge(dd^c\varphi)^k\wedge\beta^{n-k-1}.
\]
Therefore, we get
\begin{equation} \label{eq:estimationB}
|B (\delta)| \leq 2 \kappa \delta^{\alpha} \, I'_k(v,\varphi),
\end{equation}
where $I'_k (v,\varphi) := \int_{\Omega'} dd^c v \wedge(dd^c\varphi)^k\wedge\beta^{n-k-1} < + \infty$ by Chern-Levine-Nirenberg inequalities (see \cite{Lu12}).

The problem is to estimate the masses $I'_k (v,\varphi) $ by a uniform constant which does not depend on $\Omega' \Subset \Omega$.
We could use the obvious inequality
\[
\int_{\Omega'} dd^c v \wedge(dd^c\varphi)^k\wedge\beta^{n-k-1} \leq \int_{\Omega} dd^c v \wedge(dd^c\varphi)^k\wedge\beta^{n-k-1},
\]
to conclude under the finiteness full mass condition
\[ I_k (v,\varphi) := \int_{\Omega} dd^c v \wedge(dd^c\varphi)^k\wedge\beta^{n-k-1}< + \infty.\]
This is precisely the statement of Lemma \ref{lem:claim2} when $g \in C^{1,1} (\partial \Omega)$, whose justification was missing in the proof of \cite[Lem.\,4.2]{BZ20}.

So let us proceed to the proof the second statement of the lemma assuming that $g \in C^{1,1} (\partial \Omega)$.
By \eqref{eq:estimationB} and the formula \eqref{eq:claim2} of Lemma \ref{lem:claim2} applied to $v$ and $\psi = \varphi$, we get
\begin{equation} \label{eq:estimationC}
|B (\delta)| \leq \kappa \, d(m,n) \, R \, \delta^{\alpha},
\end{equation}
where $d(m,n) = d(m,n,\varphi,g) >0$ is a uniform constant.

For simplicity set $ \sigma_{k }(\varphi) := (dd^c \varphi)^k \wedge \beta^{n - k}.$ Combining the inequalities \eqref{eq:estimationA}, \eqref{eq:estimationB} and \eqref{eq:estimationC}, we obtain for $\delta$ such that $0 < \delta < \delta_0$,
\begin{equation} \label{eq:funestimate}
\int_{\Omega}(u-v) \sigma_{k +1}(\varphi)\leq M_1 \frac{\kappa \delta^\alpha}{\delta^2} \int_{\Omega} (u-v) \sigma_{k }(\varphi) + \kappa \, d(m,n) \, R \, \delta^{\alpha}.
\end{equation}

To finish the proof of the last statement of the lemma, we argue by induction on~$k$ for $k$ such that $0 \leq k \leq m$.
When $k=0$, the inequality is obviously satisfied with $C_0 = 1$ and $\alpha_0 = 1$.

Assume that the inequality holds for some integer $0 \leq k\leq m - 1$, \ie
\begin{equation}\label{eq:Hyprec}
\int_{\Omega}( u-v) \sigma_k{(\varphi)}\leq C_{k} R[\Vert u-v \Vert_1]^{\alpha_k}.
\end{equation}
We will show that there exists $C_{k+1} > 0$ such that
\[
\int_{\Omega}(u-v) \sigma_{k+1}(\varphi)\leq C_{k+1} R[\Vert u-v \Vert_1]^{\alpha_{k+1}}.
\]
Indeed \eqref{eq:funestimate} and \eqref{eq:Hyprec} yields
\[
\int_{\Omega}(u-v) \sigma_{k+1}(\varphi)\leq
M_1 C_k R \kappa \, \delta^{\alpha - 2} [\|u-v\|_1]^{\alpha_k} + d(m,n) \kappa R \delta^{\alpha}.
\]

We want to optimize the last estimate by a suitable choice of $\delta$. Since $ \|u-v\|_1 \leq 1$, we can take $\delta= \delta_0 [\|u-v\|_1]^{\alpha_k\slash 2} < \delta_0$ in the last inequality to obtain
\begin{align*}
\int_{\Omega}(u-v) \sigma_{k+1}(\varphi) &\leq  (M_1 C_k
+ d (m,n))\, \kappa \, R \, \bigl(\|u-v\|_1]^{\alpha_k \slash 2} \bigr)^{\alpha } \\
& \leq C_{k + 1} R [\|u-v\|_1]^{\alpha_{k+ 1}},
\end{align*}
where $\alpha_{k + 1} := \alpha_k (\alpha \slash 2)$.
This proves the last statement of the lemma.

We now proceed to the proof of the first statement.
As we saw before the main issue is to estimate uniformly the integrals like $ I'_k (v,\varphi)$.
We do not know whether the integrals $I_k (v,\varphi)$ are finite. We will rather estimate the following integrals which behave much better:
\[
J_{k} (u,v,\varphi) := \int_{\Omega} (u-v)^m (dd^c \varphi)^{k} \wedge \beta^{n -k }.
\]
By Hölder inequality, we have
\begin{equation} \label{eq:Hineq2}
\int_\Omega (u-v) (dd^c \varphi)^k \wedge \beta^{n-k} \leq \vert \Omega\vert^{ (m-1) \slash m}\left(J_k (u,v,\varphi)\right)^{1 \slash m},
\end{equation}
where $\vert \Omega\vert$ is the volume of $\Omega$.

It is then enough to estimate $J_k (u,v,\varphi)$. We will proceed by induction on~$k$ \hbox{($0 \leq k \leq m$)} to prove the following estimate
\begin{equation} \label{eq:Jk}
J_k (u,v) \leq C'_k R (1 + \Vert u - v\Vert_{\infty}^{m -1} ) \, \|u-v\|_1^{\alpha_k},
\end{equation}
where $C'_k = C'_k (k,m,n\varphi,g) > 0$ is a constant.

If $k = 0$, we have $J_0 (u,v) \leq \Vert u - v\Vert_{\infty}^{m - 1} \Vert u - v\Vert_1$. Assume the estimate \eqref{eq:Jk} is proved for some integer $0 \leq k \leq m-1$. To prove it for the integer $k + 1$, we write as before for $\delta$ such that $0 < \delta < \delta_0$,
\[
J_{k+1} (u,v) = A' (\delta) + B' (\delta),
\]
where
\[
A' (\delta):=\int_{\Omega}(u-v)^m dd^c \varphi_{\delta}\wedge(dd^c\varphi)^{k} \wedge\beta^{n-k},
\]
and
\[
B' (\delta):= \int_{\Omega}(u-v)^m dd^c(\varphi - \varphi_{\delta} - \kappa \delta^\alpha)\wedge(dd^c\varphi)^{k} \wedge\beta^{n-k}.
\]
The first term $A' (\delta)$ is estimated as before using \eqref{eq:Jk}, \ie
\begin{equation} \label{eq:estimationA'}
|A' (\delta)|\leq M_1 \,\frac{\kappa \delta^{\alpha}}{\delta^2}\, J_k (u,v).
\end{equation}
We need to estimate the second term $B' (\delta)$. Since $ u - v = 0$ near the boundary, we~can integrate by parts to get the following formula
\begin{equation} \label{eq:estimationB'0}
B' (\delta) = \int_{\Omega} (\varphi_{\delta} - \varphi + \kappa \delta^\alpha) \left(- dd^c[ (u-v)^m]\right) \wedge (dd^c\varphi)^k\wedge\beta^{n-k-1}.
\end{equation}
If $m = 1$, then $k = 0$. Since $dd^c u \wedge \beta^{n-1} \geq 0$ weakly on $\Omega$, it follows that
\[
- dd^c [(u-v) \wedge \beta^{n-1} \leq dd^c v \wedge \beta^{n-1},
\]
weakly on $\Omega$.
Since $0 \leq \varphi_{\delta} - \varphi + \kappa \delta^\alpha \leq 2 \kappa \delta^\alpha$, it follows from \eqref{eq:estimationB'0} that
\begin{equation*} \label{eq:estimateB'1}
\vert B' (\delta) \vert \leq 2 \kappa \delta^\alpha \int_\Omega dd^c v \wedge\beta^{n-1} \leq 2 R \, \kappa \, \delta^\alpha.
\end{equation*}
If $ m \geq 2$, a simple computation shows that
\begin{align*}
- dd^c [(u-v)^m] & ={} - m (u-v)^{m -1} dd^c (u - v) \\
&\hphantom{{}={}}- m (m-1) (u-v)^{m - 2} d(u-v) \wedge d^c(u-v) \\
&\leq{}  - m (u-v)^{m -1} dd^c (u- v),
\end{align*}
Since $ dd^c u \wedge (dd^c\varphi)^k\wedge\beta^{n-k-1} \geq 0$ weakly on $\Omega$, it follows that
\[
- dd^c [(u-v)^m] \wedge (dd^c\varphi)^k\wedge\beta^{n-k-1} \leq m (u-v)^{m -1} dd^c v\wedge (dd^c\varphi)^k\wedge\beta^{n-k-1},
\]
weakly on $\Omega$.
Since $0 \leq \varphi_{\delta} - \varphi + \kappa \delta^\alpha \leq 2 \kappa \delta^\alpha$, it follows from \eqref{eq:estimationB'0} that
\begin{equation} \label{eq:estimation2}
|B' (\delta)| \leq 2 m \, \kappa \, \delta^{\alpha} \int_\Omega (u-v)^{m -1} dd^c v \wedge (dd^c\varphi)^k\wedge\beta^{n-k-1}.
\end{equation}

Recall that $\beta = dd^c \psi_0$, where $\psi_0 (z) := \vert z\vert^2 -r_0^2$, where $r_0 > 0$ is chosen so that $\psi_0 \leq 0$ on $\Omega$.
Then the inequality \eqref{eq:estimation2} implies that
\begin{equation} \label{eq:estimationB'}
|B' (\delta)| \leq 2 m \kappa \delta^{\alpha} \int_\Omega (u-v)^{m -1} dd^c v \wedge (dd^c\varphi)^k \wedge (dd^c \psi)^{m - k-1} \wedge\beta^{n-m}.
\end{equation}
Since $v, \varphi, \psi_0 \leq 0$ on $\Omega$, repeating the integration by parts $(m-1)$ times, we deduce from \eqref{eq:estimationB'} that
\begin{equation} \label{eq:estimationC'}
|B' (\delta)| \leq 2 m! \, \kappa \, \delta^{\alpha} \Vert \varphi\Vert_{\infty}^k \Vert \psi_0\Vert_{\infty}^{m - k-1} \int_\Omega (dd^c v)^m \wedge\beta^{n-m}.
\end{equation}

Combining the inequalities \eqref{eq:estimationA'} and \eqref{eq:estimationC'}, we obtain for $\delta$ such that $0 < \delta < \delta_0$,
\[
\int_{\Omega}(u-v)^m \sigma_{k +1}(\varphi)\leq M_1 \kappa \delta^{\alpha- 2} J_k (u,v)+ d'(m,n) R \delta^{\alpha},
\]
where $d'(m,n) = d'(m,n,\varphi,g) > 0$ is a uniform constant.
Applying the induction hypothesis \eqref{eq:Jk}, we get
\[
\int_{\Omega}(u-v)^m \sigma_{k +1}(\varphi)\leq \kappa M_1 C_k\delta^{\alpha- 2}  (1+\Vert u - v \Vert_{\infty}^{m -1}) \|u-v\|_1^{\alpha_k}+ d'(m,n) R \delta^{\alpha}.
\]
We want to optimize the last estimate. Since $ \|u-v\|_1 \leq 1$, we can take
\[
\delta= \delta_0 [\|u-v\|_1]^{\alpha_k\slash 2} < \delta_0
\]
in the last inequality to obtain
\begin{align*}
\int_{\Omega}(u-v)^m \sigma_{k+1}(\varphi) &\leq  \bigl(\kappa M_1 C_k (1+\Vert u - v\Vert_{\infty}^{m -1}) + R d'(m,n)\bigr)\, \bigl(\|u-v\|_1^{\alpha_k \slash 2}\bigr)^{\alpha } \\
& \leq  C_{k + 1} (1 + \Vert u - v\Vert_{\infty}^{m -1} ) R\, \|u-v\|_1^{\alpha_{k+ 1}},
\end{align*}
where $\alpha_{k + 1} := \alpha_k (\alpha \slash 2)$ and and $C_{k+1} := \kappa M_1 \delta_0^{\alpha-2}C_k + d'(m,n)$.
This proves the estimate \eqref{eq:Jk} for $k+1$. Taking into account the inequality \eqref{eq:Hineq2} we obtain the estimate of the lemma with appropriate constants. This finishes the proof of the second part of the lemma.
\end{proof}

\section{Proof of Theorem \ref{th:B'}}

We are now ready to prove Theorem \ref{th:B'} using Lemma \ref{lem:ModC}.

\begin{proof} The strategy of the proof is the same as that of \cite[Th.\,B]{BZ20}, but for convenience of the reader we will recall the main steps.

We know that the Dirichlet problem \eqref{eq:DirPb} admits a unique bounded solution $u \in \mathcal{SH}_m (\Omega) \cap L^{\infty} (\Omega)$ by \cite{N13}.

To prove that $u$ is Hölder continuous with some exponent $0 < \theta < 1$, it is enough to consider the standard regularization $u_\delta := u \star \chi_\delta$, for $\delta$ such that $0 < \delta < \delta_0$ and to prove that there exists a constant $L_1 > 0$ and $0< \delta_1 < \delta_0$ such that for any $\delta$ such that $0 < \delta < \delta_1$, and any $z \in \Omega_\delta$, we have
\begin{equation} \label{eq:Holdercontinous}
\sup_{\Omega_\delta} (u_\delta - u) \leq C \delta^{\theta},
\end{equation}
provided we prove that $u$ is Hölder continuous near the boundary with the same exponent $\theta$. We refer the reader to \cite{Ze20} for a complete proof of this fact.

Recall that $u_\delta$ can be defined for $\delta$ such that $0 < \delta \leq \delta_0$ (where $\delta_0 >0$ is small enough) by the following formula : for $z \in \C^n$,
\[
u_\delta (z) = \int_{\Omega} u (\zeta) \chi_\delta (z-\zeta) d \lambda_{2 n} (\zeta),
\]
where $\lambda_{2 n}$ is the Lebesgue measure on $\Omega$.
Observe that $u_\delta $ is smooth on $\C^n$ and $m$-subharmonic on $\Omega_\delta$.

The first step is to use the Hölder continuity of the boundary datum $g$ to construct global barriers to show that $u$ is Hölder continuous near the boundary and to deduce bounded $m$-subharmonic global approximants $(\tilde u_\delta)_{0 < \delta < \delta_0}$ of $u$ on $\Omega$ which satisfy the following inequalities:
\begin{equation} \label{eq:bestapproximant}
0 \leq \tilde{u}_{\delta}(z) - u (z) \leq u_\delta (z) - u (z) \leq \tilde{u}_{\delta}(z) - u (z) + \kappa \delta^\alpha, \quad z\in \Omega_{\delta},
\end{equation}
and $ \tilde{u}_{\delta} \geq u$ on $ \Omega$ and $\tilde{u}_{\delta} = u$ on $\Omega \setminus \Omega_\delta$ for $\delta$ such that $0 < \delta < \delta_0$.
We refer to \cite{BZ20} for this construction.

Since $\tilde u_\delta = u$ near the boundary $\partial \Omega$, it follows from Stokes formula that for $\delta$ such that $0 < \delta < \delta_0$, we have
\[
\int_{\Omega}(dd^c\tilde{u}_{\delta})^m\wedge\beta^{n-m} = \int_{\Omega}(dd^c u)^m\wedge\beta^{n-m} \leq \mu (\Omega) < \infty.
\]

The second step is to use stability estimates. Since $\tilde u_\delta = u$ on $\Omega \setminus \Omega_\delta$ and $\mu\leq (dd^c\varphi)^m\wedge\beta^{n-m}$ on $\Omega$, we can apply \cite[Prop.\,2.20]{BZ20} with the exponent $\tau$ given in \cite[Th.\,A]{BZ20}, to deduce that for any $\gamma$ such that
\[
0 <\gamma < \gamma' (m,n,\alpha):=
\frac{m \alpha}{ m (m + 1) \alpha + (n-m) [(2 - \alpha)m + \alpha]},
\]
there exists a constant $D_\gamma > 0$ such that any $0 < \delta < \delta_0$,
\begin{equation} \label{eq2}
\sup_{\Omega}(\tilde{u}_{\delta}-u) \leq D_{\gamma} \biggl(\int_{\Omega}(\tilde{u}_{\delta}-u) d\mu\biggr)^{\gamma}.
\end{equation}

To prove \eqref{eq:Holdercontinous}, we need to estimate $\int_{\Omega}(\tilde{u}_{\delta}-u) d\mu$ in terms of the integral \hbox{$\int_{\Omega}(\tilde{u}_{\delta}-u) d\lambda_{2n}$}, where $\lambda_{2n}$ is the Lebesgue measure.

Since $\mu \leq (dd^c \varphi)^m \wedge \beta^{n-m}$, this is possible by applying Lemma \ref{lem:ModC} if we can ensure that $\Vert (\tilde{u}_{\delta}-u \Vert_{1} := \int_{\Omega}(\tilde{u}_{\delta}-u) d\lambda_{2n} \leq 1$ for $\delta > 0$ small enough. This is clearly true since $\tilde{u}_\delta$ decreases to $ u$ as $\delta$ decreases to $0$, but we will need a quantitative estimate.

Indeed by \cite{GKZ08} (see also \cite[Lem.\,2.3]{BZ20}), we know that
\begin{equation} \label{eq3}
\int_{\Omega}({u}_{\delta}-u) d \lambda_{2 n} \leq B \|\Delta u\|_{\Omega_\delta}\delta^2.
\end{equation}

To complete the proof of our theorem, it suffices to estimate the mass $\|\Delta u\|_{\Omega_\delta} := \int_{\Omega_\delta} dd^c u \wedge \beta^{n-1}$ in terms of $\delta$.
We will consider the two cases separately.

\begin{enumerate}
\item
Assume first that $g \in C^{1,1} (\partial \Omega)$.
By the inequality \eqref{eq:claim2} of Lemma \ref{lem:claim2}, we have for~$\delta$ such that $0 < \delta < \delta_0$,
\[
\|\Delta u\|_{\Omega_\delta} \leq \int_\Omega dd^c u \wedge \beta^{n - 1}\leq d(m,n) \bigl(M' + \mu (\Omega)^{1 \slash m}\bigr) < + \infty.
\]
Hence by \eqref{eq3} it follows that
\begin{equation*} \label{eq4}
\Vert {u}_{\delta} - u \Vert_1 = \int_{\Omega}({u}_{\delta}-u)d\lambda_{2n} \leq C'_m \delta^2 \leq 1,
\end{equation*}
for $\delta$ such that $0 < \delta \leq \delta_1$ with $0 < \delta_1 < \delta_0$ and small enough. Therefore, applying the inequality \eqref{eq:MocEst1} of Lemma \ref{lem:ModC}, we get for $\delta$ such that $0 < \delta < \delta_1$,
\begin{align*}
\int_{\Omega}(\tilde{u}_{\delta}-u)d\mu & \leq C_m \biggl(\int_{\Omega}(\tilde{u}_{\delta}-u)(z)d\lambda_{2 n} (z)\biggr)^{\alpha_m} \\
&\leq C_m\biggl(\int_{\Omega_{\delta}}({u}_{\delta}(z)-u(z)d\lambda_{2 n} (z)\biggr)^{\alpha_m},
\end{align*}
where the last inequality follows from \eqref{eq:bestapproximant}. Using \eqref{eq3} we deduce that for $\delta$ such that $0 < \delta < \delta_1$,
\begin{equation} \label{eq:3}
\int_{\Omega}(\tilde{u}_{\delta}-u)d\mu \leq C_m' \delta^{2 \alpha_m}.
\end{equation}
It follows from \eqref{eq:bestapproximant}, \eqref{eq2} and \eqref{eq:3} that for $\delta$ such that $0 < \delta < \delta_1$,
\begin{equation*} \label{eq:finalestimate}
\begin{split}
\sup_{\Omega_{\delta}}( {u}_{\delta}-u)&\leq\sup_{\Omega}(\tilde{u}_{\delta}-u)+ \kappa \delta^{\alpha} \\
&\leq  {C''_m} \delta^{2\gamma\alpha_m} + \kappa \delta^{\alpha}.
\end{split}
\end{equation*}
This proves the required estimate \eqref{eq:Holdercontinous} with $\theta = 2 \gamma \alpha_m < \alpha$.

\item
In the general case we need to estimate $\|\Delta u\|_{\Omega_\delta}$ in a different way. Since the defining function $\rho$ of $\Omega$ is smooth and $\vert \nabla \rho \vert > 0$ on $\partial \Omega$, it follows from Hopf's lemma that there exists a uniform constant $\tilde c_1 > 0$ such that $- \rho (z) \geq \tilde c_1 \, \text{dist} (z, \partial \Omega)$ (see \cite{Ze20} for more details). Then
\[
\int_{\Omega_\delta} dd^c u \wedge \beta^{n-1} \leq \tilde c_2 \delta^{-1} \int_{\Omega} (- \rho) dd^c u \wedge \beta^{n-1}.
\]
We can assume that $u \leq 0$ on $\Omega$. Then by the integration by parts inequality (see \cite[(2.5)]{BZ20}), it follows that
\begin{equation} \label{eq:Laplace}
\int_{\Omega_\delta} dd^c u \wedge \beta^{n-1} \leq \tilde c_2 \delta^{-1} \int_{\Omega} (- u) \, dd^c \rho \wedge \beta^{n-1}\leq \tilde c_3 \delta^{-1} \, \text{osc}_{\Omega} u,
\end{equation}
where $\tilde c_2, \tilde c_3 > 0$ are uniform constants. Hence by \eqref{eq3} and \eqref{eq:Laplace} we get
\begin{equation*} \label{eq5}
\int_{\Omega}({u}_{\delta}-u)d\lambda_{2n} \leq {\tilde C}'_m \delta \leq 1,
\end{equation*}
for $\delta$ such that $0 < \delta \leq \delta_2$ with $0 < \delta_2 < \delta_1$ and small enough.

Therefore, applying the inequality \eqref{eq:MocEst2} of Lemma \ref{lem:ModC}, we get for $\delta$ such that $0 < \delta < \delta_2$,
\begin{align*}
\int_{\Omega}(\tilde{u}_{\delta}-u)d\mu & \leq \tilde C'_m \, \left(\int_{\Omega}(\tilde{u}_{\delta}-u) d\lambda_{2 n} (z)\right)^{\tilde \alpha_m} \\
&\leq \tilde C'_m \, \left(\int_{\Omega_{\delta}}({u}_{\delta}-u)d\lambda_{2 n} \right)^{\tilde \alpha_m},
\end{align*}
where the last inequality follows from \eqref{eq:bestapproximant}. This implies that for $\delta$ such that \hbox{$0\!<\!\delta\!<\!\delta_2$},
\begin{equation} \label{eq:6}
\int_{\Omega}(\tilde{u}_{\delta}-u)d\mu \leq \tilde C_m'' \, \delta^{\tilde \alpha_m}.
\end{equation}
Therefore, taking into account \eqref{eq2} and \eqref{eq:6}, we get for $\delta$ such that $0 < \delta < \delta_2$,
\[
\sup_{\Omega_{\delta}} ({u}_{\delta}-u) \leq \tilde C_m \, \delta^{\gamma\alpha_m \slash m} + \kappa \delta^{\alpha},
\]
since $\Vert \tilde u_{\delta} - u\Vert_{\infty}\!\leq\!\text{osc}_{\Omega} u$.
This proves the required estimate \eqref{eq:Holdercontinous} with \hbox{$\theta\!=\!\gamma\alpha_m/m\!<\!\alpha$}.\qedhere
\end{enumerate}
\end{proof}
